
\section{Routing overhead, lapse factor, and gravitational time delay}
\label{app:overhead_lapse}

This appendix records the basic dynamical interface used in Sections~\ref{sec:dark_sector} and \ref{sec:dynamical_closure}.

\subsection{Overhead and lapse}
Let $\kappa(x)$ denote a local routing/compilation overhead, interpreted as a readout-time cost or scheduling depth.
Fix a reference value $\kappa_0$ and define the lapse-like factor
\[
N(x):=\frac{\kappa_0}{\kappa(x)}.
\]
The operational reading is: larger overhead implies smaller lapse.

\subsection{Time dilation and redshift (effective metric form)}
In a conservative effective metric model, one may treat the lapse as determining the temporal component of the metric,
\[
g_{00}(x)\approx -N(x)^2,
\]
so that proper time and coordinate time relate as $\dd\tau \approx N\,\dd t$.
Spatial variations in $N$ then produce gravitational redshift and time-delay effects.
In particular, regions of large overhead ($N\ll 1$) exhibit slow effective clocks and increased signal delay.

\subsection{Overhead field \texorpdfstring{$\chi$}{chi}}
Define the dimensionless overhead ratio and its logarithm,
\[
s(x)=\frac{\kappa(x)}{\kappa_0}=\frac{1}{N(x)},
\qquad
\chi(x)=\log s(x).
\]
This $\chi$ field is the natural scalar proxy used in the action-based closure and in the data reconstruction protocol (Appendix~\ref{app:data_protocol}).



